\chapter{Non-Euclidean geometry}\label{geometry:ch:noneuclid}

\begin{center}\itshape
I have created a new, another world out of nothing.
\upshape(J\'anos Bolyai, letter to his father, 3 November 1823)
\end{center}

Euclid's \emph{Elements} rest on five postulates. Four of them are short and obvious. The fifth is long: if a line falling on two lines makes the interior
angles on one side less than two right angles, the two lines, if produced indefinitely, meet on that side. It is equivalent to the statement that through a
point outside a line there is exactly one parallel. For two thousand years mathematicians tried to deduce it from the other four. In 1733 the Jesuit
Girolamo Saccheri published \emph{Euclid freed of every flaw}, in which he assumed the fifth postulate false and drew consequence after consequence,
hoping for a contradiction; he declared one found, but the ``contradiction'' was only a strangeness. A century later three men took the strangeness
seriously. Nikolai Lobachevsky in Kazan published a geometry with many parallels in 1829. J\'anos Bolyai, a young Hungarian officer, wrote to his father
the line of our epigraph and published his work as an appendix to his father's textbook in 1831--32. When the father sent it to Gauss, the answer came
back: to praise it would amount to praising himself, for he had had the same ideas long before. The new geometry was consistent, but nobody could
\emph{see} it until Eugenio Beltrami (1868) and Henri Poincaré (1882) built models of it inside Euclidean geometry.

In this chapter we build those models and prove the main facts: the hyperbolic distance, its isometries (the M\"obius maps of Appendix~C), the
parallels, and the beautiful formula for the area of a triangle, $\pi$ minus the sum of its angles. We start on the sphere, where the same formula appears
with the opposite sign.

\begin{think}
On the Earth, walk $1000$ km south from the North Pole, turn left by $90^\circ$, walk $1000$ km, turn left by $90^\circ$ again and walk back to the pole.
What is the angle at the pole between your outgoing and incoming paths? What is the sum of the three angles of your triangle?
\end{think}

\section{Spherical geometry}

On the unit sphere $S^2\subset\mathbb R^3$, the shortest paths are arcs of \emph{great circles}, the intersections of $S^2$ with planes through $0$ (a classical fact that we will
not need: here a triangle is simply defined by arcs of great circles), and the distance between $u,v\in S^2$ is the angle
$\arccos(u\cdot v)$. Two great circles always meet, in two antipodal points: the planes through $0$ that cut them out share a line. So the sphere is a
geometry without parallels. A \emph{spherical triangle} is bounded by three arcs of great circles; its angles are the angles between the arcs.

\begin{theorem}[Girard, 1629]\label{geometry:thm:girard}
A spherical triangle with angles $\alpha,\beta,\gamma$ has area $\alpha+\beta+\gamma-\pi$.
\end{theorem}

\begin{proof}
Two great circles meeting at an angle $\theta$ cut the sphere into four \emph{lunes}, two of angle $\theta$ and two of angle $\pi-\theta$. Rotating about the
axis through the two meeting points shows that the area of a lune is proportional to its angle, and the whole sphere, of area $4\pi$, is a lune of angle
$2\pi$; so a lune of angle $\theta$ has area $2\theta$.

Let $T$ be the triangle, with vertices $A,B,C$, and $T'$ the antipodal triangle $-T$, which has the same area. The two great circles through $A$ bound a lune of
angle $\alpha$ containing $T$, and its antipodal lune contains $T'$. Do the same at $B$ and $C$. These six lunes cover the sphere; points of $T$ and $T'$ are
covered three times and all other points once (except for boundaries, which have area zero). Hence
\[
2(2\alpha+2\beta+2\gamma)=4\pi+2\cdot2\operatorname{area}(T),
\]
which gives $\operatorname{area}(T)=\alpha+\beta+\gamma-\pi$.
\end{proof}

\begin{figure}[ht]
\centering
\fig{geometry-ch10-sphere}
\caption{The spherical triangle with vertices $e_1,e_2,e_3$ has three right angles. Its area is $\frac{3\pi}2-\pi=\frac\pi2$, one eighth of the sphere.}
\label{geometry:fig:sphere}
\end{figure}

So the angle sum of a spherical triangle exceeds $\pi$, by an amount equal to its area. Small triangles are almost Euclidean, large ones are not. The same
is visible in the trigonometry.

\begin{proposition}[spherical law of cosines]\label{geometry:prop:sphcos}
In a spherical triangle with sides $a,b,c$ (opposite to $A,B,C$) and angle $\gamma$ at $C$,
\[
\cos c=\cos a\cos b+\sin a\sin b\cos\gamma.
\]
\end{proposition}

\begin{proof}
Think of $A,B,C$ as unit vectors, so $\cos a=B\cdot C$, $\cos b=A\cdot C$, $\cos c=A\cdot B$. The unit tangent vectors at $C$ pointing along the arcs to $A$
and to $B$ are $t_A=\frac{A-(A\cdot C)C}{\sin b}$ and $t_B=\frac{B-(B\cdot C)C}{\sin a}$ (the component of $A$ orthogonal to $C$ has length $\sin b$). Their
inner product is $\cos\gamma$:
\[
\cos\gamma=t_A\cdot t_B=\frac{A\cdot B-(A\cdot C)(B\cdot C)}{\sin a\sin b}=\frac{\cos c-\cos a\cos b}{\sin a\sin b}.
\]
\end{proof}

For a right angle, $\cos c=\cos a\cos b$, the spherical Pythagorean theorem; for small sides, $\cos x\approx1-\frac{x^2}2$ turns it into $c^2\approx a^2+b^2$
(Exercise~\ref{geometry:ex:sphpythag}).

\section{The hyperbolic plane}

The sphere has constant curvature $+1$ (Appendix~L). We want a surface of constant curvature $-1$. Hilbert proved in 1901 that no complete smooth surface
in $\mathbb R^3$ will do, so we build it abstractly, on the upper half-plane
\[
\mathbb H=\{z=x+iy\in\mathbb C:y>0\},
\]
by declaring that a small Euclidean displacement $dz$ at height $y$ has hyperbolic length $\frac{|dz|}y$. Near the real axis, distances are stretched
enormously: the real axis is infinitely far away.

\begin{definition}
The \emph{hyperbolic length} of a piecewise $C^1$ path $\gamma:[0,1]\to\mathbb H$ is $L(\gamma)=\int_0^1\frac{|\gamma'(t)|}{\operatorname{Im}\gamma(t)}\,dt$.
The \emph{hyperbolic distance} $d(z,w)$ is the infimum of $L(\gamma)$ over the paths from $z$ to $w$. The \emph{hyperbolic area} of a region $R$ is
$\iint_R\frac{dx\,dy}{y^2}$.
\end{definition}

Angles are the Euclidean ones: the hyperbolic length element is the Euclidean one multiplied by a positive factor, which changes no angle.

\begin{proposition}[isometries]\label{geometry:prop:isometries}
Let $g(z)=\frac{az+b}{cz+d}$ with $a,b,c,d$ real and $ad-bc=1$. Then $g$ maps $\mathbb H$ onto itself and preserves hyperbolic lengths, distances and areas.
So does the reflection $z\mapsto-\bar z$.
\end{proposition}

\begin{proof}
A direct computation gives
\[
\begin{gathered}
\operatorname{Im}g(z)=\frac{\operatorname{Im}\bigl((az+b)(c\bar z+d)\bigr)}{|cz+d|^2}=\frac{(ad-bc)\operatorname{Im}z}{|cz+d|^2}=\frac{\operatorname{Im}z}{|cz+d|^2},\\
g'(z)=\frac{ad-bc}{(cz+d)^2}=\frac1{(cz+d)^2}.
\end{gathered}
\]
So $g$ maps $\mathbb H$ to $\mathbb H$, and its inverse, of the same form, maps $\mathbb H$ back: $g$ is a bijection of $\mathbb H$. For a path $\gamma$,
$|(g\circ\gamma)'|=|g'(\gamma)||\gamma'|=\frac{|\gamma'|}{|c\gamma+d|^2}$ and $\operatorname{Im}g(\gamma)=\frac{\operatorname{Im}\gamma}{|c\gamma+d|^2}$, so the
integrands of $L(g\circ\gamma)$ and $L(\gamma)$ agree. Lengths of paths are preserved, hence distances. For areas, the Jacobian of the holomorphic map $g$ is
$|g'|^2=|cz+d|^{-4}$, which exactly compensates the factor $\operatorname{Im}(g(z))^{-2}=|cz+d|^4\operatorname{Im}(z)^{-2}$. The reflection $z\mapsto-\bar z$
keeps $y$ and $|dz|$, and its Jacobian has absolute value $1$.
\end{proof}

The real matrices of determinant $1$ form the group $SL_2(\mathbb R)$, and $g$ and $-g$ give the same map. Tome~8 studies this group and its action on
$\mathbb H$ in depth. For us, the point is that $\mathbb H$ is very symmetric. The \emph{hyperbolic lines}, or \emph{h-lines}, are the vertical half-lines and
the half-circles centred on the real axis: the generalised circles orthogonal to $\mathbb R$, intersected with $\mathbb H$.

\begin{lemma}\label{geometry:lem:hlines}
Through two distinct points $z,w\in\mathbb H$ passes exactly one h-line, and there is an isometry $g$ as in Proposition~\ref{geometry:prop:isometries} with $g(z)=i$
and $g(w)=it$ for some $t>1$; it maps the h-line through $z$ and $w$ onto the imaginary axis.
\end{lemma}

\begin{proof}
If $\operatorname{Re}z=\operatorname{Re}w$, the vertical line through them is an h-line, and no half-circle centred on $\mathbb R$ contains two points with the
same real part. Otherwise the centre of a half-circle through $z$ and $w$ must lie on the perpendicular bisector of the segment $zw$, which is not vertical
and so meets $\mathbb R$ in exactly one point $c$: the half-circle centred at $c$ through $z$ is the unique h-line. Let $a<b$ be its endpoints on
$\mathbb R$ (for a vertical h-line $x=a$, take $b=\infty$). The map $g_1(z)=\frac{z-a}{b-z}$, whose matrix $\begin{pmatrix}1&-a\\-1&b\end{pmatrix}$ has
determinant $b-a>0$ (divide by $\sqrt{b-a}$; for $b=\infty$ use $g_1(z)=z-a$), sends $a$ to $0$ and $b$ to $\infty$. By Theorem~\ref{geometry:thm:mobius} it maps the h-line,
a generalised circle through $a$ and $b$ orthogonal to $\mathbb R$, to a generalised circle through $0$ and $\infty$ orthogonal to $\mathbb R$: the imaginary axis.
So $g_1(z)=is$ and $g_1(w)=is'$ with $s,s'>0$. The scaling $z\mapsto z/s$, with matrix $\operatorname{diag}(s^{-1/2},s^{1/2})$, sends them to $i$ and $it$; if
$t<1$, compose with $z\mapsto-1/z$, which fixes $i$ and sends $it$ to $i/t$.
\end{proof}

\begin{theorem}[hyperbolic distance]\label{geometry:thm:hypdist}
For $z,w\in\mathbb H$,
\[
\cosh d(z,w)=1+\frac{|z-w|^2}{2\operatorname{Im}z\operatorname{Im}w}.
\]
The arc of h-line from $z$ to $w$ is the unique shortest path between them.
\end{theorem}

\begin{proof}
First let $z=i$ and $w=it$ with $t>1$. A path $\gamma=x+iy$ from $i$ to $it$ has
\[
L(\gamma)=\int_0^1\frac{\sqrt{x'^2+y'^2}}y\,dt\ge\int_0^1\frac{|y'|}y\,dt\ge\int_0^1\frac{y'}y\,dt=\log t,
\]
with equality if and only if $x'=0$ and $y'\ge0$ throughout: the path runs up the imaginary axis. So $d(i,it)=\log t$, attained only by the segment of
h-line. The right side of the formula is $1+\frac{(t-1)^2}{2t}=\frac{t^2+1}{2t}=\cosh\log t$, so the formula holds.

In general, take $g$ from Lemma~\ref{geometry:lem:hlines}. It preserves $d$ by Proposition~\ref{geometry:prop:isometries}, and it preserves the right side of the formula too:
$\operatorname{Im}g(z)=\frac{\operatorname{Im}z}{|cz+d|^2}$ and $g(z)-g(w)=\frac{z-w}{(cz+d)(cw+d)}$ (a two-line computation using $ad-bc=1$), so the factors
$|cz+d|^2|cw+d|^2$ cancel. Shortest paths are preserved too, and $g$ maps h-lines to h-lines.
\end{proof}

\begin{figure}[ht]
\centering
\fig{geometry-ch10-halfplane}
\caption{Left: the h-line $x^2+y^2=1$ (blue) and three h-lines through $P$ that never meet it (red). Right: the region above the unit circle between
$x=\cos(\pi-\alpha)$ and $x=\cos\beta$ is a triangle with angles $\alpha$, $\beta$ and $0$ (its third vertex is at infinity).}
\label{geometry:fig:halfplane}
\end{figure}

\begin{proposition}[many parallels]\label{geometry:prop:parallels}
Let $\ell$ be an h-line and $P$ a point not on $\ell$. Infinitely many h-lines through $P$ do not meet $\ell$.
\end{proposition}

\begin{proof}
By Lemma~\ref{geometry:lem:hlines} (applied to two points of $\ell$) we may assume that $\ell$ is the imaginary axis, and by the reflection $z\mapsto-\bar z$ that
$P=x+iy$ with $x>0$. The half-circle centred at $c>0$ through $P$ has radius $r=|P-c|$, and it stays in $\operatorname{Re}z>0$ when $c-r>0$, that is, when
$c^2>(x-c)^2+y^2$, that is, when $c>\frac{x^2+y^2}{2x}$. Each such $c$ gives an h-line through $P$ disjoint from $\ell$; so does the vertical line
$\operatorname{Re}z=x$.
\end{proof}

\section{The area of a triangle}

A \emph{hyperbolic triangle} is bounded by three arcs of h-lines. We also allow \emph{ideal} vertices, on $\mathbb R\cup\{\infty\}$, where two sides meet at
infinity with angle $0$.

\begin{think}
Is there a hyperbolic triangle with three angles of $60^\circ$? With three angles of $0^\circ$? How large can its area be?
\end{think}

\begin{theorem}[Gauss--Bonnet for triangles]\label{geometry:thm:hyparea}
A hyperbolic triangle with angles $\alpha,\beta,\gamma$ has area $\pi-\alpha-\beta-\gamma$.
\end{theorem}

\begin{proof}
\emph{Step 1: one ideal vertex.} Suppose that $\gamma=0$, at an ideal vertex. An isometry sends this vertex to $\infty$ (if it is $a\in\mathbb R$, use
$z\mapsto-1/(z-a)$), so the two sides through it become vertical half-lines and the third side an arc of a half-circle centred on $\mathbb R$. The maps
$z\mapsto z+u$ ($u$ real) and $z\mapsto\lambda z$ ($\lambda>0$) are isometries (matrices $\begin{pmatrix}1&u\\0&1\end{pmatrix}$ and
$\operatorname{diag}(\lambda^{1/2},\lambda^{-1/2})$), so we may assume that this half-circle is the unit circle. The triangle is then the region
$x_0\le x\le x_1$, $y\ge\sqrt{1-x^2}$ of Figure~\ref{geometry:fig:halfplane}. At the vertex $e^{i\varphi}$ the angle between the vertical side and the circle equals the
angle between the horizontal and the radius, so the vertices are $e^{i(\pi-\alpha)}$ and $e^{i\beta}$: $x_0=\cos(\pi-\alpha)$ and $x_1=\cos\beta$. The area is
\[
\int_{x_0}^{x_1}\int_{\sqrt{1-x^2}}^{\infty}\frac{dy}{y^2}\,dx=\int_{x_0}^{x_1}\frac{dx}{\sqrt{1-x^2}}=\arccos x_0-\arccos x_1=(\pi-\alpha)-\beta.
\]
This also covers $\alpha=0$ or $\beta=0$ (then $x_0=-1$ or $x_1=1$).

\emph{Step 2: the general case.} Let $ABC$ have finite vertices, and extend the side from $A$ through $C$ beyond $C$ until it reaches the boundary at an ideal
point $D$. The triangle $ABD$ is the union of $ABC$ and $CBD$, which overlap only along $BC$. Let $\delta$ be the angle of $CBD$ at $B$. By Step 1,
$\operatorname{area}(ABD)=\pi-\alpha-(\beta+\delta)$ and, since the angle of $CBD$ at $C$ is $\pi-\gamma$, $\operatorname{area}(CBD)=\pi-(\pi-\gamma)-\delta=
\gamma-\delta$. Subtracting, $\operatorname{area}(ABC)=\pi-\alpha-\beta-\gamma$. A triangle with two finite vertices and one ideal vertex is Step 1, and one
with ideal vertices only is too.
\end{proof}

\begin{corollary}\label{geometry:cor:anglesum}
The angles of a hyperbolic triangle add up to less than $\pi$, and its area is less than $\pi$. A triangle with three ideal vertices has area exactly $\pi$.
\end{corollary}

So the hyperbolic plane is the mirror image of the sphere: angle sums fall short of $\pi$ by the area. Since the area of a triangle is determined by its
angles, there are no similar triangles of different sizes, and there is a natural unit of length: the curvature. Exercise~\ref{geometry:ex:circles} shows another
sign of negative curvature: the circumference of a circle of radius $r$ is $2\pi\sinh r$, which grows exponentially.

\section{The Poincar\'e disk}

The half-plane puts one boundary point, $\infty$, in a special position. The \emph{Cayley map} $C(z)=\frac{z-i}{z+i}$ sends $\mathbb H$ onto the open unit disk
$\mathbb D$ (a point is closer to $i$ than to $-i$ exactly when it lies in $\mathbb H$), and $i$ to the centre $0$.

\begin{proposition}[the disk model]\label{geometry:prop:disk}
Transported by $C$, the hyperbolic length element becomes $\frac{2|dw|}{1-|w|^2}$, and the distance becomes
\[
\cosh d(w_1,w_2)=1+\frac{2|w_1-w_2|^2}{(1-|w_1|^2)(1-|w_2|^2)}.
\]
The h-lines become the diameters and the arcs of circles orthogonal to the unit circle. In particular $d(0,w)=2\operatorname{artanh}|w|$.
\end{proposition}

\begin{proof}
For $w=C(z)$, $1-|w|^2=\frac{|z+i|^2-|z-i|^2}{|z+i|^2}=\frac{4y}{|z+i|^2}$ and $C'(z)=\frac{2i}{(z+i)^2}$, so $|dw|=\frac{2|dz|}{|z+i|^2}$ and
$\frac{2|dw|}{1-|w|^2}=\frac{|dz|}y$. Similarly $C(z_1)-C(z_2)=\frac{2i(z_1-z_2)}{(z_1+i)(z_2+i)}$, so
\[
\frac{2|w_1-w_2|^2}{(1-|w_1|^2)(1-|w_2|^2)}=\frac{8|z_1-z_2|^2}{16\,y_1y_2}=\frac{|z_1-z_2|^2}{2y_1y_2},
\]
and Theorem~\ref{geometry:thm:hypdist} gives the formula. The h-lines of $\mathbb H$ are the generalised circles orthogonal to $\mathbb R$; $C$ maps $\mathbb R$ to
the unit circle and preserves generalised circles and angles (Theorem~\ref{geometry:thm:mobius}). For $w_1=0$, $\cosh d=\frac{1+|w|^2}{1-|w|^2}$, which is
equivalent to $|w|=\tanh\frac d2$.
\end{proof}

In the disk, the isometries fixing $0$ include the rotations $w\mapsto e^{i\theta}w$ and the reflection $w\mapsto\bar w$, and the reflection in an h-line is
the inversion in its circle (Exercise~\ref{geometry:ex:inversioniso}). Reflecting a triangle with angles $\frac\pi2,\frac\pi3,\frac\pi7$ in its sides, again and
again, tiles the whole disk (Figure~\ref{geometry:fig:tiling}): around the vertex with angle $\frac\pi7$, fourteen triangles fit. The area of each one is
$\pi-\frac\pi2-\frac\pi3-\frac\pi7=\frac\pi{42}$, the smallest area of any triangle whose reflections tile the plane in this way. In 1958 H.~S.~M.~Coxeter sent
the artist M.~C.~Escher a paper containing such a tiling; Escher replied that the figure gave him ``quite a shock'', and it led to his woodcuts \emph{Circle
Limit I} to \emph{IV}.

\begin{figure}[ht]
\centering
\fig{geometry-ch10-tiling}
\caption{The Poincar\'e disk tiled by hyperbolic triangles with angles $\frac\pi2$, $\frac\pi3$, $\frac\pi7$, alternately coloured. All the triangles are
congruent: they only look smaller near the boundary, which is infinitely far away.}
\label{geometry:fig:tiling}
\end{figure}

\begin{remark}
The claim that $\frac\pi{42}$ is the smallest area is a classical fact about reflection groups (the triangle with angles $\frac\pi p,\frac\pi q,\frac\pi r$ tiles
when $p,q,r$ are integers, and $1-\frac1p-\frac1q-\frac1r$ is smallest for $(2,3,7)$, see Exercise~\ref{geometry:ex:tilings}). The group generated by the three
reflections, and its index-$2$ subgroup of rotations, play a starring role in the theory of Riemann surfaces: a famous theorem of Hurwitz (1893) bounds
the number of symmetries of a compact Riemann surface of genus $g\ge2$ by $84(g-1)$, and $84=2\pi/\frac\pi{42}$ comes from this triangle.
\end{remark}

\section{Exercises}

\begin{exercise}[$\star$]\label{geometry:ex:lune}
Use Theorem~\ref{geometry:thm:girard} to find the area of a spherical triangle with three right angles, and of an equilateral spherical triangle with angles of
$120^\circ$. How many of each tile the sphere?
\end{exercise}

\begin{exercise}[$\star$]\label{geometry:ex:sphpythag}
Deduce from Proposition~\ref{geometry:prop:sphcos} that $\cos c=\cos a\cos b$ when $\gamma=\frac\pi2$, and that $c^2=a^2+b^2+O(|a|^4+|b|^4)$ for small $a,b$.
\end{exercise}

\begin{exercise}[$\star$]\label{geometry:ex:sphcircle}
Show that on the unit sphere the circle of points at distance $r<\pi$ from the North Pole has length $2\pi\sin r$, less than $2\pi r$.
\end{exercise}

\begin{exercise}[$\star$]\label{geometry:ex:hpdist}
Find the h-line through $1+i$ and $3+i$, and the hyperbolic distance between them. Compare with the length $\int_1^3dx=2$ of the horizontal segment.
\end{exercise}

\begin{exercise}[$\star$]\label{geometry:ex:inversioniso}
Show that $z\mapsto1/\bar z$, the inversion in the unit circle, is an isometry of $\mathbb H$ that fixes the h-line $|z|=1$ pointwise. Deduce that the
inversion in any half-circle centred on $\mathbb R$ is an isometry.
\end{exercise}

\begin{exercise}[$\star\star$]\label{geometry:ex:circles}
Show that the hyperbolic circle of radius $r$ centred at $0$ in the disk is the Euclidean circle of radius $\tanh\frac r2$, that its length is $2\pi\sinh r$, and
that the disk it bounds has area $4\pi\sinh^2\frac r2$.
\end{exercise}

\begin{exercise}[$\star\star$]\label{geometry:ex:hypcos}
Prove the hyperbolic law of cosines: in a triangle with sides $a,b,c$ and angle $\gamma$ opposite $c$,
$\cosh c=\cosh a\cosh b-\sinh a\sinh b\cos\gamma$. Deduce a hyperbolic Pythagorean theorem.
\end{exercise}

\begin{exercise}[$\star$]\label{geometry:ex:ideal}
Show that any two triangles with three ideal vertices are congruent (there is an isometry mapping one onto the other).
\end{exercise}

\begin{exercise}[$\star\star$]\label{geometry:ex:parallelangle}
(Lobachevsky's angle of parallelism) Let $\ell$ be the imaginary axis, $F=i$, and $P=e^{i\theta}$ with $0<\theta<\frac\pi2$, so that the unit circle is the
h-line through $P$ perpendicular to $\ell$ at $F$. Let $d=d(P,F)$. Show that the h-line through $P$ and the endpoint $0$ of $\ell$ makes the angle $\theta$
with $PF$ at $P$, and that $\tan\frac\theta2=e^{-d}$.
\end{exercise}

\begin{exercise}[$\star$]\label{geometry:ex:tilings}
Let $p,q,r\ge2$ be integers such that a hyperbolic triangle with angles $\frac\pi p,\frac\pi q,\frac\pi r$ exists. Show that
$\frac1p+\frac1q+\frac1r<1$, and that its area is at least $\frac\pi{42}$, with equality only for $\{2,3,7\}$.
\end{exercise}

\begin{exercise}[$\star\star$]\label{geometry:ex:quadrilateral}
Show that a convex hyperbolic quadrilateral has angle sum less than $2\pi$; deduce that there are no rectangles in the hyperbolic plane.
\end{exercise}

\begin{exercise}[$\star\star\star$]\label{geometry:ex:horocycle}
The horizontal line $y=1$ is not an h-line. Show that it is the set of points ``at the same distance from $\infty$'': for every $z$ with $\operatorname{Im}z=1$,
$d(z,iT)-\log T\to0$ as $T\to\infty$. (Such curves are called \emph{horocycles}.) Show that the arc of $y=1$ from $i$ to $x+i$ has length $|x|$, while
$d(i,x+i)=2\operatorname{arsinh}\frac{|x|}2$.
\end{exercise}

\section{Hints}

\begin{hint}{geometry:ex:lune}
The tiling by octants comes from the three coordinate planes.
\end{hint}
\begin{hint}{geometry:ex:sphpythag}
Expand $\cos x=1-\frac{x^2}2+O(x^4)$.
\end{hint}
\begin{hint}{geometry:ex:sphcircle}
The points at distance $r$ from $(0,0,1)$ have height $\cos r$.
\end{hint}
\begin{hint}{geometry:ex:hpdist}
The centre is on the perpendicular bisector $x=2$.
\end{hint}
\begin{hint}{geometry:ex:inversioniso}
Write $1/\bar z$ as the composition of $z\mapsto-1/z$ and $z\mapsto-\bar z$.
\end{hint}
\begin{hint}{geometry:ex:circles}
Use Proposition~\ref{geometry:prop:disk} with $w_1=0$, and integrate $\frac{4\,dA}{(1-|w|^2)^2}$ in polar coordinates.
\end{hint}
\begin{hint}{geometry:ex:hypcos}
Put $C$ at $0$ in the disk, $A$ on the positive real axis and $B$ at angle $\gamma$.
\end{hint}
\begin{hint}{geometry:ex:ideal}
Use Theorem~\ref{geometry:thm:mobius}(2) with real points.
\end{hint}
\begin{hint}{geometry:ex:parallelangle}
The h-line through $P$ and $0$ is a half-circle centred at $c$ with $|P-c|=c$. Look at the isosceles triangle $0,c,P$.
\end{hint}
\begin{hint}{geometry:ex:tilings}
The condition is $\frac1p+\frac1q+\frac1r<1$. Minimise $1-\frac1p-\frac1q-\frac1r$ subject to it.
\end{hint}
\begin{hint}{geometry:ex:quadrilateral}
Cut along a diagonal.
\end{hint}
\begin{hint}{geometry:ex:horocycle}
Use Theorem~\ref{geometry:thm:hypdist} and $\operatorname{arcosh}u=\log(u+\sqrt{u^2-1})$.
\end{hint}

\section{Solutions}

\begin{solution}{geometry:ex:lune}
\idea Area is angle excess.
\writeup Three right angles: area $\frac{3\pi}2-\pi=\frac\pi2$, and $\frac{4\pi}{\pi/2}=8$ such triangles tile the sphere (the octants). Angles $\frac{2\pi}3$:
area $2\pi-\pi=\pi$, and $4$ tile the sphere: project a regular tetrahedron inscribed in the sphere from its centre; at each vertex three faces meet with
angles $\frac{2\pi}3$.
\end{solution}

\begin{solution}{geometry:ex:sphpythag}
\idea The spherical laws become the Euclidean ones to second order.
\writeup With $\cos\gamma=0$ the law reads $\cos c=\cos a\cos b$. Expanding, $1-\frac{c^2}2+O(c^4)=\bigl(1-\frac{a^2}2\bigr)\bigl(1-\frac{b^2}2\bigr)+O(a^4+b^4)=
1-\frac{a^2+b^2}2+O(a^4+b^4)$, since $a^2b^2\le\frac{a^4+b^4}2$. Also $c=O(|a|+|b|)$ (triangle inequality), so $c^2=a^2+b^2+O(a^4+b^4)$.
\end{solution}

\begin{solution}{geometry:ex:sphcircle}
\idea A circle of latitude.
\writeup A point at distance $r$ from $N=(0,0,1)$ has $u\cdot N=\cos r$: height $\cos r$, so it lies on the horizontal circle of Euclidean radius $\sin r$,
whose length is $2\pi\sin r<2\pi r$ for $r>0$.
\end{solution}

\begin{solution}{geometry:ex:hpdist}
\idea Centre on the perpendicular bisector, then the distance formula.
\writeup The centre $c$ is on $x=2$ and on $\mathbb R$: $c=2$, and the radius is $|1+i-2|=\sqrt2$. The h-line is the half-circle $(x-2)^2+y^2=2$. By
Theorem~\ref{geometry:thm:hypdist}, $\cosh d=1+\frac{4}{2}=3$, so $d=\operatorname{arcosh}3=\log(3+2\sqrt2)\approx1.763$, shorter than the horizontal segment, of
length $2$: the shortest path bulges upwards, where lengths are cheaper.
\end{solution}

\begin{solution}{geometry:ex:inversioniso}
\idea Compose two known isometries.
\writeup The map $z\mapsto-1/z$ has matrix $\begin{pmatrix}0&-1\\1&0\end{pmatrix}$ of determinant $1$, and $z\mapsto-\bar z$ is an isometry. Their composition is
$z\mapsto-\overline{(-1/z)}=1/\bar z$, an isometry. For $|z|=1$, $1/\bar z=z$, so the h-line $|z|=1$ is fixed pointwise. The inversion in the half-circle of centre
$c$ and radius $r$ is $z\mapsto c+\frac{r^2}{\overline{z-c}}$, the conjugate of $z\mapsto1/\bar z$ by the isometry $z\mapsto c+rz$; so it is an isometry.
\end{solution}

\begin{solution}{geometry:ex:circles}
\idea Centre the circle at $0$, where the disk model is rotationally symmetric.
\writeup By Proposition~\ref{geometry:prop:disk}, $d(0,w)=r$ if and only if $|w|=\rho=\tanh\frac r2$: a Euclidean circle. Its hyperbolic length is
$2\pi\rho\cdot\frac2{1-\rho^2}=2\pi\sinh r$, because $\frac{2\tanh(r/2)}{1-\tanh^2(r/2)}=\sinh r$. The area element of the disk model is $\frac{4\,dA}{(1-|w|^2)^2}$
(the square of the length factor), so the area is
\[
\int_0^\rho\frac{4\cdot2\pi s\,ds}{(1-s^2)^2}=\frac{4\pi}{1-\rho^2}-4\pi=\frac{4\pi\rho^2}{1-\rho^2}=4\pi\sinh^2\frac r2,
\]
using $\frac{\tanh^2x}{1-\tanh^2x}=\sinh^2x$. For small $r$ these are $\approx2\pi r$ and $\approx\pi r^2$; for large $r$ they grow like $e^r$.
\end{solution}

\begin{solution}{geometry:ex:hypcos}
\idea Put the angle $\gamma$ at the centre of the disk, where everything is explicit.
\writeup By isometries, $C=0$, $A=\rho_b=\tanh\frac b2$ on the positive real axis and $B=\rho_ae^{i\gamma}$ with $\rho_a=\tanh\frac a2$ (h-lines through $0$ are
diameters, and angles are Euclidean). By Proposition~\ref{geometry:prop:disk}, with $|A-B|^2=\rho_a^2+\rho_b^2-2\rho_a\rho_b\cos\gamma$,
\[
\cosh c=1+\frac{2(\rho_a^2+\rho_b^2-2\rho_a\rho_b\cos\gamma)}{(1-\rho_a^2)(1-\rho_b^2)}=\frac{(1+\rho_a^2)(1+\rho_b^2)-4\rho_a\rho_b\cos\gamma}{(1-\rho_a^2)(1-\rho_b^2)}.
\]
Since $\frac{1+\rho^2}{1-\rho^2}=\cosh x$ and $\frac{2\rho}{1-\rho^2}=\sinh x$ for $\rho=\tanh\frac x2$, this is $\cosh a\cosh b-\sinh a\sinh b\cos\gamma$. For
$\gamma=\frac\pi2$: $\cosh c=\cosh a\cosh b$, and expanding $\cosh x=1+\frac{x^2}2+\dots$ gives back $c^2\approx a^2+b^2$ for small triangles.
\end{solution}

\begin{solution}{geometry:ex:ideal}
\idea Ideal triangles are determined by three boundary points, and M\"obius maps move any three real points to any three others.
\writeup An ideal triangle is determined by its three vertices $x_1,x_2,x_3\in\mathbb R\cup\{\infty\}$. By Theorem~\ref{geometry:thm:mobius}(2) there is a M\"obius map
$g$ with $g(x_j)=0,1,\infty$. It has real coefficients (the explicit formula in the proof uses only the $x_j$), so it is $z\mapsto\frac{az+b}{cz+d}$ with
$a,b,c,d$ real and $ad-bc\ne0$. If $ad-bc>0$, divide by $\sqrt{ad-bc}$: $g$ is an isometry. If $ad-bc<0$, the map $z\mapsto-g(z)=\frac{-az-b}{cz+d}$ has determinant $-(ad-bc)>0$, so after normalising it is an isometry, and it sends the $x_j$ to
$0,-1,\infty$; the reflection $z\mapsto-\bar z$ then sends these to $0,1,\infty$. In both cases an isometry maps the triangle onto the one with vertices $0,1,\infty$.
\end{solution}

\begin{solution}{geometry:ex:parallelangle}
\idea An isosceles triangle hides in the picture.
\writeup The h-line through $P$ ending at $0$ is a half-circle centred at a real $c$ with $|P-c|=|0-c|=c$: $1-2c\cos\theta+c^2=c^2$, so $c=\frac1{2\cos\theta}$.
The triangle with vertices $0$, $c$, $P$ is isosceles with $|cP|=|c0|$, so its angles at $0$ and at $P$ are equal; the angle at $0$ is $\theta=\arg P$. The angle
between two circles at a common point is the angle between their radii there; the radii at $P$ are $P-0$ and $P-c$, and the angle between them is the angle of the
triangle at $P$, namely $\theta$. So the asymptotic parallel makes the angle $\theta$ with the perpendicular $PF$. For the distance,
$\cosh d=1+\frac{|e^{i\theta}-i|^2}{2\sin\theta}=1+\frac{2-2\sin\theta}{2\sin\theta}=\frac1{\sin\theta}$. With $u=\tan\frac\theta2$, $\sin\theta=\frac{2u}{1+u^2}$, so
$\cosh d=\frac{u^{-1}+u}2$, which means $e^{d}=u^{-1}$ (the other root, $e^d=u$, is excluded because $u<1\le e^d$). Hence $\tan\frac\theta2=e^{-d}$. As $d\to0$
the angle tends to $\frac\pi2$, as in Euclid; as $d\to\infty$ it tends to $0$.
\end{solution}

\begin{solution}{geometry:ex:tilings}
\idea The angle sum must fall short of $\pi$.
\writeup By Corollary~\ref{geometry:cor:anglesum} we need $\frac1p+\frac1q+\frac1r<1$. (Conversely such a
triangle always exists; a numerical calculation constructs it.) Its area is $\pi\bigl(1-\frac1p-\frac1q-\frac1r\bigr)$. With
$p\le q\le r$: if $p\ge3$, then $(3,3,4)$ is the worst case, giving $1-\frac{11}{12}=\frac1{12}$; if $p=2$ and $q\ge4$, the worst is $(2,4,5)$, giving $\frac1{20}$;
if $p=2$, $q=3$, we need $r\ge7$, and $r=7$ gives $1-\frac12-\frac13-\frac17=\frac1{42}$. The smallest is $\frac1{42}$, only for $\{2,3,7\}$.
\end{solution}

\begin{solution}{geometry:ex:quadrilateral}
\idea A diagonal cuts a quadrilateral into two triangles.
\writeup A convex quadrilateral $ABCD$ is cut by the diagonal $AC$ into two triangles, whose angles add up to the angles of the quadrilateral. Each triangle has
angle sum less than $\pi$, so the quadrilateral has angle sum less than $2\pi$, and its area is $2\pi$ minus its angle sum. A rectangle would have angle sum
$2\pi$.
\end{solution}

\begin{solution}{geometry:ex:horocycle}
\idea Distances to points far up the imaginary axis are almost the vertical distance.
\writeup For $z=x+i$, Theorem~\ref{geometry:thm:hypdist} gives $\cosh d(z,iT)=1+\frac{x^2+(T-1)^2}{2T}=\frac{T}2+\frac{x^2+1}{2T}$. Since $\operatorname{arcosh}u=\log\bigl(u+\sqrt{u^2-1}\bigr)=\log(2u)+O(u^{-2})$, and here $2u=T+\frac{x^2+1}T$,
\[
d(z,iT)=\log T+\log\Bigl(1+\frac{x^2+1}{T^2}\Bigr)+O(T^{-2})\longrightarrow\log T,
\]
in the sense that $d(z,iT)-\log T\to0$, independently of $x$. Along $y=1$ the length
element is $|dx|$, so the arc from $i$ to $x+i$ has length $|x|$. The distance is given by $\cosh d=1+\frac{x^2}2$, that is $\sinh\frac d2=\frac{|x|}2$
(using $\cosh d=1+2\sinh^2\frac d2$), so $d=2\operatorname{arsinh}\frac{|x|}2$, which is much smaller than $|x|$ for large $x$: horocycles are far from straight.
\end{solution}

\section*{Sources and further reading}

\begin{itemize}
\item The MacTutor biography of J\'anos Bolyai (University of St Andrews): the letter of 3 November 1823 (epigraph), the Appendix sent to Gauss in 1831, Gauss's
reply ``to praise it would amount to praising myself'', and Bolyai's discovery in 1848 that Lobachevsky had published similar work in 1829.
\item G.~Saccheri, \emph{Euclides ab omni naevo vindicatus}, Milan, 1733; E.~Beltrami, \emph{Saggio di interpretazione della geometria non-euclidea}, 1868;
H.~Poincar\'e, \emph{Th\'eorie des groupes fuchsiens}, Acta Mathematica 1 (1882).
\item A.~Girard, \emph{Invention nouvelle en l'alg\`ebre}, Amsterdam, 1629.
\item D.~Dunham, \emph{Hyperbolic art and the poster design}, Mathematics Awareness Month 2003 essay (Coxeter's 1957 paper \emph{Crystal symmetry and its
generalizations}, sent to Escher in 1958, and Escher's ``quite a shock'').
\item J.~W.~Anderson, \emph{Hyperbolic Geometry}, Springer, 2nd ed., 2005; A.~F.~Beardon, \emph{The Geometry of Discrete Groups}, Springer, 1983 (the half-plane and
disk models, triangle groups).
\item J.~Stillwell, \emph{Geometry of Surfaces}, Springer, 1992; M.~J.~Greenberg, \emph{Euclidean and Non-Euclidean Geometries}, Freeman, 4th ed., 2008 (the history
of the parallel postulate).
\end{itemize}
